% 390_Puzzleteile_Etabliert_De_ch.tex
% Johann Pascher, 9. Okt. 2026
% Lose Puzzleteile: Kaluza-Klein und andere etablierte Bausteine der Zeit-Masse-Dualität

\providecommand{\BM}{\textbf{[B]}}
\providecommand{\KM}{\textbf{[K]}}
\providecommand{\SM}{\textbf{[S]}}
\providecommand{\XM}{\textbf{[X]}}
\providecommand{\QM}{\textbf{[Q]}}
\providecommand{\EM}{\textbf{[E]}}
\providecommand{\SetzM}{\textbf{[SETZUNG]}}

\chapter*{Lose Puzzleteile: Kaluza-Klein und andere etablierte Bausteine der Zeit-Masse-Dualität}
\addcontentsline{toc}{chapter}{Lose Puzzleteile: Kaluza-Klein und andere etablierte Bausteine}

\begin{abstract}
Wer $\tilde T\cdot m=1$ zum ersten Mal sieht, erkennt darin oft etwas Bekanntes, am häufigsten
Kaluza-Klein. Das ist kein Einwand, sondern ein Befund: Fast jeder Baustein der FFGFT liegt seit
Jahrzehnten in der Literatur, allerdings verstreut über Relativitätstheorie,
Quantenmechanik, Gruppentheorie, Stringtheorie, Zahlentheorie und Phänomenologie. Dieses Dokument
ordnet elf solcher Bausteine ein: Kaluza-Klein, de Broglies innere Uhr und die Compton-Zeit, die
Zitterbewegung, Paulis Einwand gegen einen Zeitoperator, die relationalen Zeitansätze, Wigners
Masse als Casimir-Invariante, den Casimir-Wert der $SU(3)$, die Heaviside-Lorentz-Einheiten, die
Orbifold-Kompaktifizierung, die endlichen Körper mit ihrem Frobenius-Automorphismus und die
Koide-Formel. Zu jedem Baustein wird festgehalten, was er bereits enthielt, welches Stück fehlte
und an welcher Stelle des Korpus er sich einfügt. Wichtig ist die Reihenfolge: Diese Bausteine
waren nicht die Grundlage der FFGFT. Die Theorie entstand unabhängig von ihnen aus der Beziehung
$\tilde T\cdot m=1$; die Bausteine fanden sich erst nachträglich in der Literatur und fügen sich
in ein Bild, das ohne sie gezeichnet wurde. Das Ergebnis ist eine Karte der Stellen, an denen die
Theorie bereits hätte gefunden werden können.
\end{abstract}

% ============================================================
\section*{Statusmarkierungen}
% ============================================================
\begin{center}
\begin{tabular}{@{}lp{11cm}@{}}
\textbf{[SETZUNG]} & Motivierte Eingabe, keine Ableitung. \\[2pt]
\textbf{[B]} & Algebraisch bewiesen (Prüfskript, exakte Arithmetik). \\[2pt]
\textbf{[K]} & Numerisch bestätigt. \\[2pt]
\textbf{[S]} & Strukturell plausibel, Identifikation offen. \\[2pt]
\textbf{[E]} & Etablierte Mathematik oder Physik, übernommen. \\[2pt]
\textbf{[Q]} & Aus der Literatur übernommen, mit Quelle. \\[2pt]
\textbf{[X]} & Geprüft und negativ. \\
\end{tabular}
\end{center}
Die Marker in diesem Dokument geben den Status \emph{im Korpus} wieder; das Dokument selbst leitet
nichts neu her, sondern ordnet ein. Prüfskript: \texttt{pruef\_390\_puzzleteile.py}.

% ============================================================
\section{Die Frage}
% ============================================================

Im einführenden Gespräch auf der Website fragt der skeptische Physiker gleich zu Beginn, ob das
nicht nach Kaluza-Klein klinge. Die Antwort ist ja, und die Frage lässt sich verallgemeinern: Wenn
die Dualität von Masse und Zeit schon bekannt war, warum wurde sie nicht zu Ende gedacht?

Für Einstein ist diese Frage im Korpus bereits beantwortet. Dok.~312, Kapitel \enquote{Warum
Einstein $\tilde T\cdot m=1$ nicht sah}, zeigt, dass die Zutaten auf seinem Tisch lagen und dass
der fehlende Schritt nicht eine Formel war, sondern der Statuswechsel einer bekannten Formel zum
Grundsatz. Das vorliegende Dokument wiederholt diese Darstellung nicht. Es weitet den Blick: Nicht
nur ein Bestandteil, sondern fast alle lagen bereit, jeder in einem anderen Fach, jeder als
Werkzeug oder Kuriosität behandelt. Die folgenden Abschnitte gehen sie in drei Gruppen durch: die
eingerollte Richtung, die Masse als Struktur, die Geometrie und ihre Algebra.

\textbf{Zur Reihenfolge.} Keiner der folgenden Bausteine war Ausgangspunkt der FFGFT. Die Theorie
wurde aus $\tilde T\cdot m=1$ und der Geometrie des Trägers entwickelt; die Literaturstellen sind
erst nachträglich gesucht und gefunden worden. Wenn es unten heißt, ein Baustein füge sich an einer
Stelle des Korpus ein, ist das als Übereinstimmung im Nachhinein zu lesen, nicht als Herkunft. Gerade
darin liegt der Wert der Karte: Teile, die unabhängig voneinander und unabhängig von der FFGFT
entstanden sind, passen in dasselbe Bild.

% ============================================================
\section{Zeit und Masse: die eingerollte Richtung}
% ============================================================

\subsection{Kaluza-Klein (1921, 1926)}

\textbf{Was es enthielt.} Kaluza erweiterte die Allgemeine Relativitätstheorie um eine fünfte
Dimension, Klein rollte sie zu einem Kreis vom Radius $R$ ein \QM~\cite{Kaluza1921,Klein1926}. Auf
einem Kreis sind nur ganzzahlige Wicklungen zulässig, daraus entsteht ein Modenturm mit
$m_n\propto n/R$, und der Impuls entlang der Kreisrichtung erscheint in vier Dimensionen als
Masse. Klein erhielt so zugleich die Quantisierung der Ladung. Damit stand 1926 die Kernaussage,
dass eine kompakte Richtung Masse als Kehrwert ihrer Periode erzeugt.

\textbf{Was fehlte.} Die eingerollte Richtung war räumlich und kam zur Raumzeit hinzu. Der Radius
$R$ blieb ein freier Parameter, und der vorhergesagte Turm schwerer Anregungen wurde nie gefunden.
Man hatte, wie Dok.~311 es formuliert, eine Skala gewonnen und einen Parameter dazugekauft.

\textbf{Wo es sich im Korpus einfügt.} $\tilde T\cdot m=1$ ist im Korpus ausdrücklich die
Kaluza-Klein-Relation der vierten Torusrichtung, der Massendimension $S^1_m$ (Dok.~330, Dok.~270).
Die Richtung ist zeitartig: Was dort für eine eingerollte Raumrichtung gilt, gilt hier für die
eingerollte Zeit, mit der Masse in der Rolle des Modenindex (Dok.~307; ebenso Dok.~231). Der
Radius ist kein freier Modul, weil $\tilde T\cdot m=1$ die Richtung bindet; Dok.~311 fasst das als
\enquote{Skala ohne Modul} zusammen, Dok.~318 stellt die drei Wege nebeneinander. Einen
zusätzlichen Turm, auf den zu warten wäre, gibt es nicht: Das diskrete Spektrum der kompakten
Richtung ist bereits das beobachtete Teilchenspektrum (Dok.~307). Lokal gilt die
Lorentz-Invarianz uneingeschränkt, die Kompaktheit zeigt sich nur global (Dok.~312).

\subsection{De Broglies innere Uhr und die Compton-Zeit (1924, 2013)}

\textbf{Was es enthielt.} De Broglie schrieb jedem Teilchen eine innere Uhr mit der Frequenz
$\nu=mc^2/h$ zu \QM~\cite{deBroglie1925}. In natürlichen Einheiten ist das $T\cdot m=1$ in Rohform.
Die Compton-Zeit $\hbar/(mc^2)$ ist seither ein Standardbegriff, und Lan, Müller und Mitarbeiter
haben 2013 eine Uhr gebaut, deren Takt direkt aus der Masse eines Atoms folgt
\QM~\cite{Lan2013}.

\textbf{Was fehlte.} Aus der Uhr wurde die Wellenlänge; die Uhr selbst fiel als abgeleitete Skala
heraus. Die Geschichte dieses Schritts steht in Dok.~312 und wird hier nicht wiederholt.

\textbf{Wo es sich im Korpus einfügt.} Die de-Broglie-Relation erscheint auf der Massenrichtung als
$\lambda_4\cdot m=2\pi$, modenweise exakt (Dok.~314); der Uhrcharakter jeder Masse ist der
Ausgangspunkt der Theorie, nicht ihre Folgerung (Dok.~365).

\subsection{Die Zitterbewegung (1930, 1990)}

\textbf{Was es enthielt.} Schrödinger fand in der Dirac-Gleichung eine schnelle Oszillation des
Elektronenorts mit der Kreisfrequenz $2mc^2/\hbar$ \QM~\cite{Schroedinger1930}. Hestenes deutete
sie als innere Kreisbewegung und nannte das Elektron ausdrücklich eine Uhr
\QM~\cite{Hestenes1990}.

\textbf{Was fehlte.} Die Zitterbewegung blieb eine Eigenheit der Einteilchen-Dirac-Theorie, die im
Feldbild als Interferenz positiver und negativer Energieanteile gilt; eine geometrische Richtung,
in der sie stattfindet, wurde ihr nicht zugeordnet.

\textbf{Einordnung.} Im Korpus wurde die Zitterbewegung bisher nicht herangezogen. Sie gehört in
diese Reihe, weil sie dieselbe Reziprozität von Masse und Takt in einer anderen Sprache zeigt. Eine
Identifikation mit der Kreisrichtung $S^1_m$ wird hier nicht behauptet; der Faktor~2 der Frequenz
und die Rolle der negativen Energieanteile müssten dafür erst geklärt werden \SM.

\subsection{Paulis Einwand gegen einen Zeitoperator (1933)}

\textbf{Was es enthielt.} Pauli zeigte, dass zu einem nach unten beschränkten Hamiltonoperator
kein selbstadjungierter Zeitoperator existieren kann \QM~\cite{Pauli1933}. Das galt jahrzehntelang
als Grund, die Zeit aus dem Zustandsraum herauszuhalten.

\textbf{Was fehlte.} Der Einwand setzt eine offene Zeit voraus. Eine kompakte Zeit verhält sich wie
eine Winkelvariable; ihr konjugiertes Spektrum ist diskret, wie beim Drehimpuls zum Winkel. Diese
Lücke im Argument wurde nicht konsequent genutzt.

\textbf{Wo es sich im Korpus einfügt.} Dok.~307 und Dok.~334 halten fest, dass Paulis Einwand die
eingerollte Sicht nicht trifft: Für eine kompakte Zeit ist die Pauli-Bedingung nicht erfüllt, und
die Diskretheit der Moden wird zum Inhalt statt zum Hindernis.

\subsection{Relationale und zeitlose Ansätze (1983 ff.)}

\textbf{Was es enthielt.} Page und Wootters gewinnen die Zeit aus Korrelationen in einem statischen
Gesamtzustand \QM~\cite{PageWootters1983}; die Wheeler-DeWitt-Gleichung und Rovellis relationale
Zeit gehen ähnlich vor. Kaluza-Klein-Varianten mit kompakter Zeit und die Zwei-Zeit-Physik von Bars
wurden untersucht \QM~\cite{Bars2001}.

\textbf{Was fehlte.} Die Hauptlinie nimmt die Zeit als Fundamentales \emph{heraus}. Den umgekehrten
Weg, die Zeit als kompakte Dimension \emph{hineinzunehmen} und daraus die Massenmoden zu gewinnen,
hat sie kaum verfolgt.

\textbf{Wo es sich im Korpus einfügt.} Genau diese Abgrenzung zieht Dok.~307: geteilt wird die
Ablehnung der äußeren Parameterzeit, verschieden ist der positive Vorschlag, kompakte geometrische
Zeit statt relationaler oder thermischer Zeit.

% ============================================================
\section{Masse als Struktur}
% ============================================================

\subsection{Wigner: Masse als Casimir-Invariante (1939)}

\textbf{Was es enthielt.} Wigner klassifizierte die Teilchen als irreduzible Darstellungen der
Poincaré-Gruppe; die Masse ist darin ein Casimir-Label, also bereits ein struktureller Eigenwert
und keine angeheftete Zahl \EM~\cite{Wigner1939}.

\textbf{Was fehlte.} Welche Werte dieses Label annimmt, sagt die Klassifikation nicht; die
Massenwerte bleiben Eingaben.

\textbf{Wo es sich im Korpus einfügt.} Dok.~307 nennt Wigner als die Referenz dafür, dass Masse als
Strukturgröße wohlvertraut ist; Dok.~339 verwendet die Wigner-Klassifikation als Standardschritt
bei der Trennung massiver und masseloser Sektoren. Die Massenwerte selbst kommen in der FFGFT aus
der kompakten Richtung, nicht aus der Poincaré-Gruppe.

\subsection{Der Casimir-Wert der $SU(3)$}

\textbf{Was es enthielt.} Der quadratische Casimir-Operator der Fundamentaldarstellung von $SU(N)$
hat den Wert $(N^2-1)/(2N)$, für $SU(3)$ also $4/3$ \EM~\cite{Georgi1999}. Der Wert ist Lehrbuchstoff
der Farbdynamik.

\textbf{Was fehlte.} Niemand hatte Anlass, diesen Wert mit einer Modenzahl einer kompakten
Geometrie zu verknüpfen.

\textbf{Wo es sich im Korpus einfügt.} In Dok.~324 ist $\xi=C_2(SU(3)_{\mathrm{fund}})/N_{\mathrm{Fourier}}
=(4/3)/10^4=4/30000$; die Formel ist dort \BM, die Fourier-Modenzahl \KM. Derselbe Wert folgt in
Dok.~338 aus Gruppenordnungen endlicher Körper (Abschnitt~\ref{sec:galois390}).

\subsection{Heaviside-Lorentz-Einheiten (1893)}

\textbf{Was es enthielt.} Heaviside führte rationalisierte Einheiten ein, in denen die Faktoren
$4\pi$ aus den Feldgleichungen verschwinden \QM~\cite{Heaviside1893}. Mit $\hbar=c=\varepsilon_0=1$
wird $\alpha=1$ zu einer Einheitenwahl.

\textbf{Was fehlte.} Die Einheitenwahl blieb eine Rechenbequemlichkeit; dass $\hbar$ und $c$ bloße
Umrechnungsfaktoren sind, wurde erst mit der Praxis natürlicher Einheiten und der SI-Reform 2019
selbstverständlich (Dok.~312).

\textbf{Wo es sich im Korpus einfügt.} Dok.~365 stellt die Theorie in dieser Grundform auf: $\alpha=1$
ist keine Setzung, sondern eine Einheitenwahl, und die elektromagnetischen Größen werden zu
abgeleiteten Formen von Masse und Zeit. Die bekannte Zahl $137$ entsteht erst beim Übergang ins
SI-System; ihre Übereinstimmung mit dem Messwert hängt am Anker $E_0$ (Dok.~338).

% ============================================================
\section{Geometrie und Algebra}
% ============================================================

\subsection{Orbifold-Kompaktifizierung (1985)}

\textbf{Was es enthielt.} Dixon, Harvey, Vafa und Witten kompaktifizierten die Stringtheorie auf
Orbifolds, Tori geteilt durch eine endliche Gruppe; $T^6/\mathbb{Z}_3$ war eines der ersten
Beispiele \QM~\cite{DHVW1985}. Die Fixpunkte der Gruppe tragen dort eigene Sektoren, und die
Zahl der Generationen wird zu einer topologischen Größe.

\textbf{Was fehlte.} Sechs zusätzliche räumliche Dimensionen und eine riesige Zahl möglicher
Kompaktifizierungen; die Zeit bleibt außerhalb. Dok.~311 legt die beiden Bedingungen des Korpus an
die Stringtheorie an: Die kompakten Richtungen sind vorbildlich kompakt, alle Information wird in
erhaltene Quantenzahlen umgebucht, offen bleibt, was die Richtungen tragen.

\textbf{Wo es sich im Korpus einfügt.} Der Träger ist $T^4/\mathbb{Z}_3$ auf dem $D_4$-Gitter, vier
Richtungen statt sechs zusätzlicher, mit der Zeit-Masse-Richtung darunter (Dok.~330, Dok.~314).
Die $\mathbb{Z}_3$ wirkt auf allen Impulsen ungleich null frei, im Ortsraum hat der Orbifold neun
Fixpunkte (R131 in Dok.~190). Torus, $\mathbb{Z}_3$ und $D_4$ sind Setzungen \SetzM{} (Dok.~365).

\subsection{Endliche Körper und der Frobenius-Automorphismus}
\label{sec:galois390}

\textbf{Was es enthielt.} Die endlichen Körper $GF(p^n)$, ihre zyklischen multiplikativen Gruppen
und der Frobenius-Automorphismus $x\mapsto x^p$ sind seit Galois vollständig beschrieben
\EM~\cite{Lidl1997}. Mit der Physik der Elementarteilchen wurden sie nicht in Verbindung gebracht.

\textbf{Was fehlte.} Eine physikalische Struktur, in der die Primzahl~3 und ihre Erweiterungen
natürlich auftreten.

\textbf{Wo es sich im Korpus einfügt.} Die $\mathbb{Z}_3$ des Orbifolds liefert diese Struktur. Der
Frobenius-Automorphismus von $GF(9)$ entspricht der Sektorpaarung $k\leftrightarrow-k$ \BM{}
(Dok.~336). Die Ladungsquantisierung folgt aus dem Legendre-Symbol modulo~13 \BM{} (Dok.~346),
die drei Generationen ordnen sich in die Frobenius-Orbits von $GF(27)^*$ \BM{} (Dok.~348), und die
Identität $(r_\mu/r_e)^2/\xi=|GF(9)^*|^2\cdot5^2\cdot|GF(27)|=43200$ verbindet die Leptonleiter mit
den Gruppenordnungen \KM{} (Dok.~338). Die Hyperladung der geladenen Leptonen ist dabei noch
offen (R146 in Dok.~190).

\subsection{Die Koide-Formel (1981--1983)}

\textbf{Was es enthielt.} Koide fand, dass die drei geladenen Leptonmassen die Beziehung
$Q=(m_e+m_\mu+m_\tau)/(\sqrt{m_e}+\sqrt{m_\mu}+\sqrt{m_\tau})^2=2/3$ auf etwa
$10^{-5}$ erfüllen \QM~\cite{Koide1983}. Über vierzig Jahre blieb die Formel ohne Erklärung,
eine Kuriosität der Phänomenologie.

\textbf{Was fehlte.} Ein Grund für den Wert $2/3$ und für die Phase, die die Formel in ihrer
Kosinusform verlangt.

\textbf{Wo es sich im Korpus einfügt.} Die Koide-Phase ist $\theta=2/9$ \BM, derselbe Quotient wie das
Betragsquadrat einer Ikosaeder-Projektion (Dok.~367). Die Amplitude $d/c=\sqrt2$ hat den
richtigen Wert, ist aber noch nicht aus der Geometrie hergeleitet und wird als Setzung geführt
\SM{} (Dok.~353, 367; R150 in Dok.~190). Ein weiteres Puzzleteil fügt sich hier nachträglich ein:
Foot zeigte 1994, dass $Q=2/3$ einen Winkel von $45^\circ$ zwischen dem Vektor der $\sqrt{m_k}$
und der Achse $(1,1,1)$ bedeutet \QM~\cite{Foot1994}. In der Sprache des Korpus ist das die
Gleichverteilung der Massensumme auf die $\mathbb{Z}_3$-symmetrische Mode und das Paar der
nicht-trivialen Moden \BM{} (Dok.~353). Die offene Frage heißt damit nicht mehr, warum gerade
$\sqrt2$, sondern warum sich die Massenstruktur je zur Hälfte auf diese beiden Anteile verteilt.
Das exakte Paar $(\sqrt2,\,2/9)$ ist für $m_\mu/m_e$ ausgeschlossen (Dok.~369). Die Formel ist
damit eingeordnet, aber nicht vollständig geschlossen.

% ============================================================
\section{Warum die Teile nicht zusammenkamen}
% ============================================================

Drei Gründe ziehen sich durch alle Abschnitte.

\textbf{Die Teile lagen in verschiedenen Fächern.} Kaluza-Klein gehört zur Gravitation und später
zur Stringtheorie, die innere Uhr und die Zitterbewegung zur Quantenmechanik, Wigner und der
Casimir-Wert zur Gruppentheorie, die endlichen Körper zur Zahlentheorie, die Koide-Formel zur
Phänomenologie. Wer in einem dieser Fächer arbeitet, hat selten Anlass, die anderen als Teile
desselben Bildes zu lesen.

\textbf{Jedes Teil war als Werkzeug oder Kuriosität eingeordnet.} Die Compton-Zeit galt als
Umrechnung, Heavisides Einheiten als Bequemlichkeit, die Zitterbewegung als Artefakt, die
Koide-Formel als Zufall. Der fehlende Schritt war jeweils derselbe wie bei Einstein: eine bekannte
Beziehung nicht als abgeleitete Größe, sondern als Grundsatz zu nehmen (Dok.~312).

\textbf{Zwei Weichen standen falsch.} Paulis Einwand hielt die Zeit aus dem Zustandsraum heraus,
und Kaluza-Klein rollte eine räumliche statt der zeitlichen Richtung ein. Beide Weichen wären mit
derselben Einsicht umzustellen gewesen: Eine kompakte Zeit entgeht dem Einwand, und sie liefert die
Massenmoden ohne freien Radius.

Dazu kommt ein methodischer Punkt. Dass die Teile ins Bild passen, ist eine strukturelle
Übereinstimmung, keine numerische. Je komplexer die Quantenzahlen, desto näher kommt man der Rauschgrenze, an der sich
immer eine passende Kombination findet; dort wird kein Status vergeben (Dok.~342, Dok.~343).
Schemaabhängige Größen des Standardmodells wie $\Lambda_{\mathrm{QCD}}$ sind von jedem
Herleitungsanspruch ausgenommen. Die Puzzleteile tragen also nur dort, wo sie einfache, konventionsfreie
Verhältnisse betreffen, und genau dort liegen die Stellen dieses Dokuments.

% ============================================================
\section{Übersicht}
% ============================================================

\begin{center}
\small
\begin{tabular}{@{}>{\raggedright\arraybackslash}p{3.0cm}>{\raggedright\arraybackslash}p{1.7cm}>{\raggedright\arraybackslash}p{4.1cm}>{\raggedright\arraybackslash}p{3.8cm}@{}}
\toprule
Baustein & Jahr & Was fehlte & FFGFT im Korpus \\
\midrule
Kaluza-Klein & 1921/26 & räumliche Richtung, freier Radius, unbeobachteter Turm & Dok.~330, 270, 307, 311, 318 \\
De Broglie, Compton-Zeit & 1924/2013 & Uhr nur als abgeleitete Skala & Dok.~312, 314, 365 \\
Zitterbewegung & 1930/1990 & keine geometrische Richtung & bisher nicht verwendet \SM \\
Paulis Einwand & 1933 & setzt offene Zeit voraus & Dok.~307, 334 \\
Relationale Zeit, Zwei-Zeit-Physik & 1983 ff. & Zeit herausgenommen statt hineingenommen & Dok.~307 \\
Wigner & 1939 & Massenwerte bleiben Eingaben & Dok.~307, 339 \\
Casimir-Wert $SU(3)$ & Lehrbuch & keine Verbindung zu Moden & Dok.~324, 338 \\
Heaviside-Lorentz & 1893 & nur Rechenbequemlichkeit & Dok.~365, 312 \\
Orbifold & 1985 & sechs Raumrichtungen, Zeit außen & Dok.~330, 314, 311; R131 \\
Endliche Körper, Frobenius & Galois & keine physikalische Struktur & Dok.~336, 338, 346, 348 \\
Koide, Foot & 1981--83, 1994 & kein Grund für $2/3$ und die Phase & Dok.~367, 353, 369; R150 \\
\bottomrule
\end{tabular}
\end{center}

% ============================================================
\section{Fazit}
% ============================================================

Die FFGFT ist an elf Stellen anschlussfähig an etablierte Physik und Mathematik, und an jeder
dieser Stellen hätte sie gefunden werden können. Kaluza-Klein hatte die kompakte Richtung, de
Broglie die Uhr, Pauli den Einwand, dessen Lücke die kompakte Zeit ist, Wigner die Masse als
Struktur, die Gruppentheorie den Wert $4/3$, die Stringtheorie den Orbifold, die Zahlentheorie die
endlichen Körper und Koide das Rätsel, das sie zusammen lösen helfen. Die FFGFT ist nicht aus diesen
Teilen gebaut. Sie entstand unabhängig aus der Zeit als kompakter Richtung, die die Massenmoden
trägt, und aus der Algebra, die daraus mit dem einen Parameter $\xi$ folgt. Die Bausteine fanden
sich erst nachträglich in der Literatur, und sie fügen sich ein. Das spricht nicht gegen die
Theorie, sondern für sie: Teile, die unabhängig voneinander entstanden sind, passen in ein Bild,
das ohne sie gezeichnet wurde.

\vspace{1em}
\noindent\textit{Prüfskript:} \texttt{2/python/Dok390\_Skripte/pruef\_390\_puzzleteile.py}
--- 59/59.

% ============================================================
\section*{Werkzeug und Methode}
\addcontentsline{toc}{section}{Werkzeug und Methode}
% ============================================================

Claude (Anthropic) wurde als Werkzeug eingesetzt für: Zusammenstellen der Fundstellen im Korpus,
Schreiben und Ausführen des Prüfskripts, Ausarbeiten der \LaTeX-Quellen nach den Vorgaben des
Autors. Das Prüfskript sichert jede Korpusangabe dieses Dokuments durch eine Textprüfung in der
genannten Kapiteldatei ab; schlägt eine Prüfung fehl, wird die Angabe nicht verwendet.

\begin{thebibliography}{99}

\bibitem{Kaluza1921}
Th.~Kaluza, \textit{Zum Unitätsproblem der Physik},
Sitzungsber.\ Preuß.\ Akad.\ Wiss.\ Berlin (1921) 966--972.

\bibitem{Klein1926}
O.~Klein, \textit{Quantentheorie und fünfdimensionale Relativitätstheorie},
Z.\ Phys.\ \textbf{37} (1926) 895--906.

\bibitem{deBroglie1925}
L.~de~Broglie, \textit{Recherches sur la théorie des quanta},
Ann.\ Phys.\ (Paris) 10e sér.\ \textbf{3} (1925) 22--128 (Dissertation 1924).

\bibitem{Lan2013}
S.-Y.~Lan, P.-C.~Kuan, B.~Estey, D.~English, J.~M.~Brown, M.~A.~Hohensee, H.~Müller,
\textit{A clock directly linking time to a particle's mass},
Science \textbf{339} (2013) 554--557.

\bibitem{Schroedinger1930}
E.~Schrödinger, \textit{Über die kräftefreie Bewegung in der relativistischen Quantenmechanik},
Sitzungsber.\ Preuß.\ Akad.\ Wiss., Phys.-Math.\ Kl.\ \textbf{24} (1930) 418--428.

\bibitem{Hestenes1990}
D.~Hestenes, \textit{The zitterbewegung interpretation of quantum mechanics},
Found.\ Phys.\ \textbf{20} (1990) 1213--1232.

\bibitem{Pauli1933}
W.~Pauli, \textit{Die allgemeinen Prinzipien der Wellenmechanik},
in: Handbuch der Physik, Bd.~24/1, Springer, Berlin 1933.

\bibitem{PageWootters1983}
D.~N.~Page, W.~K.~Wootters, \textit{Evolution without evolution: Dynamics described by stationary
observables}, Phys.\ Rev.\ D \textbf{27} (1983) 2885--2892.

\bibitem{Bars2001}
I.~Bars, \textit{Survey of two-time physics},
Class.\ Quantum Grav.\ \textbf{18} (2001) 3113--3130.

\bibitem{Wigner1939}
E.~Wigner, \textit{On unitary representations of the inhomogeneous Lorentz group},
Ann.\ Math.\ \textbf{40} (1939) 149--204.

\bibitem{Georgi1999}
H.~Georgi, \textit{Lie Algebras in Particle Physics}, 2.~Aufl., Westview Press, 1999.

\bibitem{Heaviside1893}
O.~Heaviside, \textit{Electromagnetic Theory}, Bd.~1, The Electrician, London, 1893.

\bibitem{DHVW1985}
L.~Dixon, J.~A.~Harvey, C.~Vafa, E.~Witten, \textit{Strings on orbifolds},
Nucl.\ Phys.\ B \textbf{261} (1985) 678--686.

\bibitem{Lidl1997}
R.~Lidl, H.~Niederreiter, \textit{Finite Fields}, 2.~Aufl., Cambridge University Press, 1997.

\bibitem{Foot1994}
R.~Foot, \textit{A note on Koide's lepton mass relation}, arXiv:hep-ph/9402242 (1994).

\bibitem{Koide1983}
Y.~Koide, \textit{A fermion-boson composite model of quarks and leptons},
Phys.\ Lett.\ B \textbf{120} (1983) 161--165.

\end{thebibliography}
